\documentclass[12pt,a4paper]{amsart}
\usepackage{amsmath,amssymb,amsthm}
\usepackage{mathtools}
\usepackage[utf8]{inputenc}
\usepackage[T1]{fontenc}
\usepackage[ngerman]{babel}
\usepackage{hyperref}
\usepackage{enumitem,booktabs,array}
\usepackage{geometry}
\geometry{margin=2.5cm}

\newtheorem{theorem}{Satz}[section]
\newtheorem{lemma}[theorem]{Lemma}
\newtheorem{corollary}[theorem]{Korollar}
\newtheorem{proposition}[theorem]{Proposition}
\newtheorem{conjecture}{Vermutung}
\theoremstyle{definition}
\newtheorem{definition}[theorem]{Definition}
\newtheorem{example}[theorem]{Beispiel}
\theoremstyle{remark}
\newtheorem{remark}[theorem]{Bemerkung}

\title{Kombinatorik und Ramsey-Theorie:\\
Abz{\"a}hlen, Extremalprobleme und die Probabilistische Methode}

\author{Michael Fuhrmann}
\address{Specialist Mathematics Project, Build 144}
\email{specialist-maths@localhost}
\date{\today}

\begin{abstract}
Diese Arbeit gibt eine umfassende Darstellung der Kombinatorik und der
Ramsey-Theorie, die von klassischen Abz\"ahlprinzipien \"uber erzeugende
Funktionen und extremale Graphentheorie bis hin zu Ramsey-Zahlen und der
Probabilistischen Methode reicht.  Wir beweisen die grundlegenden S\"atze
des Fachgebiets: das Inklusions-Exklusions-Prinzip, das Schubfachprinzip,
den Ramsey-Satz (Endlichkeit von $R(s,t)$), die exakten Werte $R(3,3)=6$ und
$R(4,4)=18$, die Erd\H{o}s-Untergrenze $R(k,k) > 2^{k/2}$ mittels der
Probabilistischen Methode, die Binomial-Oberschranke
$R(k,k) \le \binom{2k-2}{k-1} \le 4^k$, den Satz von Hales--Jewett sowie
den Tur\'an-Satz f\"ur extremale Graphen.  Das Lov\'asz-Lokale-Lemma wird mit
Anwendungen auf Ramsey-Schranken und Graphenf\"arbung vorgestellt.  Abschlie\ss{}end
werden offene Probleme diskutiert, insbesondere der unbekannte exakte Wert
$R(5,5) \in [43,48]$ sowie Erd\H{o}s-Preisaufgaben.  Alle Implementierungen
basieren auf den Modulen \texttt{combinatorics.py} und \texttt{ramsey\_numbers.py}
des Specialist-Mathematics-Projekts.
\end{abstract}

\maketitle

\tableofcontents

%% ============================================================
\section{Einleitung}
%% ============================================================

Die Kombinatorik ist jener Zweig der Mathematik, der sich mit dem Z\"ahlen,
Anordnen und Ausw\"ahlen diskreter Strukturen befasst.  Ihr Einfluss reicht von
elementarer Abz\"ahlung bis zu tiefen Strukturtheorien, die Zahlentheorie,
Wahrscheinlichkeitsrechnung und Topologie miteinander verbinden.  Im Zentrum
der modernen Kombinatorik steht die Ramsey-Theorie, die in ihrer einfachsten
Form besagt: \emph{Vollst\"andige Unordnung ist unm\"oglich}.  Jede hinreichend
gro\ss{}e kombinatorische Struktur muss eine hochorganisierte Teilstruktur
enthalten.

Das Fachgebiet beginnt mit Ramseys Satz von 1930 \cite{Ramsey1930}, der
gew\"ahrleistet, dass es f\"ur beliebige positive ganze Zahlen $s$ und $t$ ein
endliches $N$ gibt, sodass jede Zweif\"arbung der Kanten des vollst\"andigen
Graphen $K_N$ entweder ein einfarbiges $K_s$ oder ein einfarbiges $K_t$ enth\"alt.
Die kleinste solche Zahl $N$ ist die \emph{Ramsey-Zahl} $R(s,t)$.

Trotz der Einfachheit dieser Aussage ist die Berechnung von Ramsey-Zahlen
au\ss{}erordentlich schwierig.  Nur eine Handvoll exakter Werte ist bekannt.
Erd\H{o}s' ber\"uhmte Bemerkung \cite{Erdos1947} bringt die Schwierigkeit
treffend auf den Punkt: W\"urde eine au\ss{}erirdische Zivilisation die Vernichtung
androhen, bis die Menschheit $R(5,5)$ berechnet, w\"are es richtig, alle
mathematischen Kr\"afte auf dieses Problem zu konzentrieren; bei der Forderung
nach $R(6,6)$ hingegen w\"are es besser, die Au\ss{}erirdischen zu vernichten.

Diese Arbeit ist wie folgt gegliedert.  Abschnitt~2 behandelt klassische
Abz\"ahlprinzipien.  Abschnitt~3 f\"uhrt erzeugende Funktionen, Catalan- und
Stirling-Zahlen ein.  Abschnitt~4 definiert Ramsey-Zahlen und beweist
$R(3,3)=6$ und $R(4,4)=18$.  Abschnitt~5 etabliert die rekursive Oberschranke
und Erd\H{o}s' probabilistische Untergrenze.  Abschnitt~6 leitet die Binomial-
Oberschranke und die Spencer-Verfeinerung her.  Abschnitt~7 behandelt den
Hales--Jewett-Satz.  Abschnitt~8 behandelt extremale Kombinatorik einschlie\ss{}lich
des Tur\'an-Satzes.  Abschnitt~9 beweist das Lov\'asz-Lokale-Lemma und gibt
Anwendungen.  Abschnitt~10 gibt einen \"Uberblick \"uber offene Probleme.

%% ============================================================
\section{Abz\"ahlprinzipien}
%% ============================================================

\subsection{Das Schubfachprinzip}

\begin{theorem}[Schubfachprinzip]\label{satz:schubfach}
Werden $n+1$ oder mehr Objekte in $n$ Schubf\"acher verteilt, so enth\"alt
mindestens ein Schubfach mindestens zwei Objekte.  Allgemeiner gilt: Werden
$kn+1$ Objekte in $n$ Schubf\"acher verteilt, so enth\"alt ein Schubfach
mindestens $k+1$ Objekte.
\end{theorem}

\begin{proof}
Angenommen, jedes Schubfach enth\"alt h\"ochstens $k$ Objekte.  Dann ist die
Gesamtzahl der Objekte h\"ochstens $kn$, im Widerspruch zur Voraussetzung von
$kn+1$ Objekten.
\end{proof}

\begin{example}
Unter 13 Personen teilen mindestens zwei denselben Geburtsmonat.
Unter $n+1$ ganzen Zahlen sind mindestens zwei kongruent modulo $n$.
\end{example}

\subsection{Das Inklusions-Exklusions-Prinzip}

Seien $A_1, A_2, \ldots, A_n$ endliche Mengen.  Definiere
\[
  S_k \;=\; \sum_{1 \le i_1 < \cdots < i_k \le n} \bigl|A_{i_1} \cap \cdots \cap A_{i_k}\bigr|.
\]

\begin{theorem}[Inklusions-Exklusions-Prinzip]\label{satz:inklusion-exklusion}
\[
  \bigl|A_1 \cup A_2 \cup \cdots \cup A_n\bigr|
  \;=\;
  \sum_{k=1}^{n}(-1)^{k-1} S_k
  \;=\;
  S_1 - S_2 + S_3 - \cdots + (-1)^{n-1}S_n.
\]
\end{theorem}

\begin{proof}
Jedes Element $x \in A_1 \cup \cdots \cup A_n$, das genau $m$ der Mengen $A_i$
angeh\"ort, wird auf der linken Seite einmal gez\"ahlt.  Auf der rechten Seite
wird es
\[
  \sum_{k=1}^{m}(-1)^{k-1}\binom{m}{k}
  = 1 - (1-1)^m = 1
\]
mal gez\"ahlt, nach dem binomischen Lehrsatz.  Die Behauptung folgt.
\end{proof}

\begin{corollary}[Anzahl der Derangementen]
Die Anzahl der Permutationen von $\{1,\ldots,n\}$ ohne Fixpunkt ist
\[
  D_n \;=\; n!\sum_{k=0}^{n}\frac{(-1)^k}{k!}
  \;\longrightarrow\; \frac{n!}{e} \quad \text{f\"ur } n\to\infty.
\]
\end{corollary}

\subsection{Doppelz\"ahlung}

Doppelz\"ahlung (auch: Zweifaches Z\"ahlen) ist eine grundlegende Technik:
Eine Identit\"at wird bewiesen, indem eine endliche Menge $S$ auf zwei
verschiedene Weisen gez\"ahlt wird.

\begin{example}[Handschlaglemma]
In jedem Graphen $G = (V,E)$ gilt
\[
  \sum_{v \in V} \deg(v) \;=\; 2|E|.
\]
Jede Kante $\{u,v\}$ tr\"agt $1$ zu $\deg(u)$ und $1$ zu $\deg(v)$ bei,
also insgesamt $2$ zur Gradfolge.
\end{example}

%% ============================================================
\section{Erzeugende Funktionen}
%% ============================================================

\subsection{Gew\"ohnliche erzeugende Funktionen}

Die \emph{gew\"ohnliche erzeugende Funktion} (OGF) einer Folge $(a_n)_{n\ge 0}$
ist die formale Potenzreihe
\[
  A(x) \;=\; \sum_{n=0}^{\infty} a_n x^n.
\]
Der Koeffizient $[x^n]\,A(x) = a_n$ kodiert die Z\"ahlinformation.

\subsection{Exponentielle erzeugende Funktionen}

F\"ur Folgen, die von gelabelten Strukturen herr\"uhren, ist die \emph{exponentielle
erzeugende Funktion} (EGF) gegeben durch
\[
  \hat{A}(x) \;=\; \sum_{n=0}^{\infty} a_n \frac{x^n}{n!}.
\]
Die Anzahl der Permutationen von $n$ Elementen ist $n!$, mit EGF
$\hat{A}(x) = \frac{1}{1-x}$.

\subsection{Catalan-Zahlen}

\begin{definition}
Die $n$-te \emph{Catalan-Zahl} ist
\[
  C_n \;=\; \frac{1}{n+1}\binom{2n}{n}, \qquad n \ge 0.
\]
\end{definition}

Die Folge beginnt mit $1, 1, 2, 5, 14, 42, 132, \ldots$.  Catalan-Zahlen z\"ahlen
viele kombinatorische Objekte: vollst\"andige Bin\"arb\"aume mit $n+1$ Bl\"attern,
Triangulierungen eines $(n+2)$-Ecks, Dyck-Pfade der L\"ange $2n$ sowie g\"ultige
Klammerausdr\"ucke der L\"ange $2n$.

\begin{proposition}[OGF der Catalan-Zahlen]
Die OGF $C(x) = \sum_{n=0}^{\infty} C_n x^n$ erf\"ullt die Funktionalgleichung
$C(x) = 1 + x\,C(x)^2$ und ergibt
\[
  C(x) \;=\; \frac{1 - \sqrt{1-4x}}{2x}.
\]
\end{proposition}

\subsection{Stirling-Zahlen}

\begin{definition}
Die \emph{Stirling-Zahl zweiter Art} $S(n,k)$ z\"ahlt die Partitionen einer
$n$-elementigen Menge in genau $k$ nicht-leere, ungelabelte Bl\"ocke.
Die \emph{Stirling-Zahl erster Art} $\left[{n\atop k}\right]$ (vorzeichenlos)
z\"ahlt Permutationen von $n$ Elementen mit genau $k$ Zyklen.
\end{definition}

Wichtige Identit\"aten:
\[
  S(n,k) = k\,S(n-1,k) + S(n-1,k-1), \qquad
  x^n = \sum_{k=0}^{n} S(n,k)\,(x)_k,
\]
wobei $(x)_k = x(x-1)\cdots(x-k+1)$ das fallende Fakult\"at ist.

%% ============================================================
\section{Ramsey-Theorie}
%% ============================================================

\subsection{Definition und grundlegende Eigenschaften}

\begin{definition}\label{def:ramsey-zahl}
F\"ur positive ganze Zahlen $s, t \ge 2$ ist die \emph{Ramsey-Zahl} $R(s,t)$
die kleinste positive ganze Zahl $N$, sodass jede Zweif\"arbung der Kanten
von $K_N$ (rot/blau) entweder ein rotes $K_s$ oder ein blaues $K_t$ enth\"alt.
\end{definition}

\"Aquivalent ist $R(s,t)$ das kleinste $N$, sodass jeder Graph auf $N$ Knoten
entweder eine Clique der Gr\"o\ss{}e $s$ oder eine unabh\"angige Menge der Gr\"o\ss{}e
$t$ enth\"alt.

\begin{remark}
Die Symmetrie $R(s,t) = R(t,s)$ folgt unmittelbar aus dem Vertauschen der Farben
(bzw. dem Komplementgraphen).  Randbedingungen: $R(2,t) = t$ und $R(s,2) = s$
f\"ur alle $s,t \ge 2$.
\end{remark}

\subsection{Der Ramsey-Satz}

\begin{theorem}[Ramsey, 1930]\label{satz:ramsey-endlich}
F\"ur alle ganzen Zahlen $s,t \ge 2$ ist die Ramsey-Zahl $R(s,t)$ endlich.
\end{theorem}

\begin{proof}
Wir beweisen $R(s,t) \le \binom{s+t-2}{s-1}$ durch Induktion.  Die Basisfälle
$R(2,t)=t$ und $R(s,2)=s$ sind klar.  F\"ur den Induktionsschritt sei $v$ ein
Knoten in einem Graphen $G$ auf $R(s-1,t)+R(s,t-1)$ Knoten.  Die Nachbarn von
$v$ bilden $N(v)$, die Nicht-Nachbarn $\overline{N}(v)$.  Nach dem Schubfachprinzip
gilt
\[
  |N(v)| \ge R(s-1,t) \quad \text{oder} \quad |\overline{N}(v)| \ge R(s,t-1).
\]
Im ersten Fall enth\"alt der Teilgraph auf $N(v)$ ein rotes $K_{s-1}$ oder ein
blaues $K_t$; zusammen mit $v$ ergibt der erste Unterfall ein rotes $K_s$.
Analoges gilt f\"ur den zweiten Fall.  Nach Induktionsvoraussetzung:
\[
  R(s,t) \le R(s-1,t) + R(s,t-1) \le \binom{s+t-3}{s-2} + \binom{s+t-3}{s-1}
  = \binom{s+t-2}{s-1}.\qedhere
\]
\end{proof}

\subsection{\texorpdfstring{$R(3,3) = 6$}{R(3,3) = 6}}

\begin{theorem}\label{satz:r33}
$R(3,3) = 6$.
\end{theorem}

\begin{proof}
\textbf{Untergrenze $R(3,3) > 5$.}
Der Kreis $C_5$ auf 5 Knoten hat kein Dreieck ($K_3$) und keine unabh\"angige
Menge der Gr\"o\ss{}e 3 (das Komplement von $C_5$ ist wieder $C_5$, das ebenfalls
kein Dreieck hat).  Daher ist $C_5$ eine Zweif\"arbung von $K_5$ ohne einfarbiges
$K_3$, was $R(3,3) > 5$ zeigt.

\textbf{Obergrenze $R(3,3) \le 6$.}
Betrachte eine beliebige Zweif\"arbung von $K_6$.  Fixiere einen Knoten $v$.
Er hat 5 Nachbarn, jede Kante rot oder blau gef\"arbt.  Nach dem Schubfachprinzip
haben mindestens 3 Kanten von $v$ dieselbe Farbe, etwa rot.  Seien
$\{u_1, u_2, u_3\}$ drei rote Nachbarn von $v$.  Ist eine Kante $u_i u_j$ rot,
so bilden $v, u_i, u_j$ ein rotes $K_3$.  Sind alle Kanten unter
$\{u_1, u_2, u_3\}$ blau, so bilden sie ein blaues $K_3$.  In jedem Fall
existiert ein einfarbiges $K_3$.
\end{proof}

\subsection{\texorpdfstring{$R(4,4) = 18$}{R(4,4) = 18}}

\begin{theorem}\label{satz:r44}
$R(4,4) = 18$.
\end{theorem}

\begin{proof}[Beweissk\"otte]
\textbf{Obergrenze $R(4,4) \le 18$.}
Aus der rekursiven Ungleichung folgt:
$R(4,4) \le R(3,4) + R(4,3) = 9 + 9 = 18$.

\textbf{Untergrenze $R(4,4) > 17$.}
Der \emph{Paley-Graph} $P(17)$ auf 17 Knoten — zwei Knoten $i,j \in \mathbb{Z}_{17}$
sind genau dann verbunden, wenn $i-j$ ein quadratischer Rest mod 17 ist — ist
selbstkomplementär und enth\"alt weder $K_4$ noch eine unabh\"angige Menge der
Gr\"o\ss{}e 4.  Daher gilt $R(4,4) > 17$.  Der Wert $R(4,4) = 18$ wurde erstmals
von Kalbfleisch (1965) etabliert und sp\"ater computergest\"utzt best\"atigt.
\end{proof}

\subsection{Tabelle bekannter Ramsey-Zahlen}

\begin{table}[h]
\centering
\caption{Bekannte exakte Werte von $R(s,t)$ f\"ur kleine $s,t$.}
\begin{tabular}{@{}ccccccc@{}}
\toprule
$R(s,t)$ & $t=2$ & $t=3$ & $t=4$ & $t=5$ & $t=6$ & $t=7$ \\
\midrule
$s=2$ & 2 & 3 & 4 & 5 & 6 & 7 \\
$s=3$ & 3 & 6 & 9 & 14 & 18 & 23 \\
$s=4$ & 4 & 9 & 18 & 25 & -- & -- \\
$s=5$ & 5 & 14 & 25 & \textbf{?} & -- & -- \\
\bottomrule
\end{tabular}
\label{tab:ramsey}
\end{table}

%% ============================================================
\section{Ramsey-Schranken}
%% ============================================================

\subsection{Die rekursive Oberschranke}

\begin{theorem}[Rekursive Oberschranke]\label{satz:rekursiv}
F\"ur alle $s,t \ge 3$ gilt
\[
  R(s,t) \;\le\; R(s-1,t) + R(s,t-1).
\]
Sind dar\"uber hinaus beide Werte $R(s-1,t)$ und $R(s,t-1)$ gerade, so gilt
$R(s,t) \le R(s-1,t) + R(s,t-1) - 1$.
\end{theorem}

\begin{proof}
Die Ungleichung $R(s,t) \le R(s-1,t) + R(s,t-1)$ wurde im Beweis von
Satz~\ref{satz:ramsey-endlich} gezeigt.  F\"ur die Parit\"atsverbesserung:
Sind beide Schranken gerade, so ist $n := R(s-1,t) + R(s,t-1) - 1$ ungerade.
Bei einer Zweif\"arbung von $K_n$ hat jeder Knoten $v$ den Grad $n-1$, der gerade
ist.  Die Parit\"atsbedingung erzwingt, dass einer der beiden F\"alle eintritt.
\end{proof}

\subsection{Die Erd\H{o}s-Untergrenze}

Die von Erd\H{o}s \cite{Erdos1947} eingeführte probabilistische Methode
revolutionierte die Kombinatorik, indem sie die Existenz kombinatorischer
Strukturen ohne explizite Konstruktion nachwies.

\begin{theorem}[Erd\H{o}s 1947 -- Probabilistische Untergrenze]\label{satz:erdos-unter}
F\"ur alle $k \ge 3$ gilt
\[
  R(k,k) > 2^{k/2}.
\]
\end{theorem}

\begin{proof}
Sei $n = \lfloor 2^{k/2} \rfloor$.  F\"ärbe jede Kante von $K_n$ unabh\"angig
und gleichverteilt rot oder blau.  F\"ur eine feste $k$-elementige Teilmenge
$S \subseteq V$ sei $A_S$ das Ereignis, dass $S$ ein einfarbiges $K_k$ bildet.
Dann gilt
\[
  \Pr[A_S] = 2 \cdot \frac{1}{2^{\binom{k}{2}}} = 2^{1-\binom{k}{2}}.
\]
Es gibt $\binom{n}{k}$ solcher Teilmengen, daher liefert die Vereinigungsschranke:
\[
  \Pr\Bigl[\bigcup_S A_S\Bigr]
  \;\le\; \binom{n}{k} \cdot 2^{1-\binom{k}{2}}
  \;\le\; \frac{n^k}{k!} \cdot 2^{1-k(k-1)/2}.
\]
F\"ur $n = \lfloor 2^{k/2}\rfloor$ gilt $n^k \le 2^{k^2/2}$.  Da
$\binom{k}{2} = k(k-1)/2$, ergibt sich
\[
  \frac{2^{k^2/2}}{k!} \cdot 2^{1-k^2/2+k/2}
  = \frac{2^{k/2+1}}{k!} < 1
\]
f\"ur alle $k \ge 3$ (da $k! > 2^{k/2+1}$ f\"ur $k \ge 3$).  Mit positiver
Wahrscheinlichkeit existiert also kein einfarbiges $K_k$, was
$R(k,k) > n \ge 2^{k/2}$ beweist.
\end{proof}

\begin{remark}
Dies ist die erste und vielleicht eindrucksvollste Anwendung der Probabilistischen
Methode.  Das Argument zeigt nur die \emph{Existenz}; der Beweis liefert keinen
expliziten Graphen auf $\lfloor 2^{k/2}\rfloor$ Knoten ohne einfarbiges $K_k$.
F\"ur kleine F\"alle liefern Paley-Graphen explizite Konstruktionen.
\end{remark}

%% ============================================================
\section{Oberschranken f\"ur diagonale Ramsey-Zahlen}
%% ============================================================

\subsection{Die Binomialschranke}

\begin{theorem}[Binomial-Oberschranke]\label{satz:binomial-oben}
F\"ur alle $k \ge 2$ gilt
\[
  R(k,k) \;\le\; \binom{2k-2}{k-1} \;\le\; 4^{k-1}.
\]
\end{theorem}

\begin{proof}
Nach Satz~\ref{satz:rekursiv} gilt $R(k,k) \le 2R(k-1,k)$.  Durch Aufl\"osen
der Rekursion und Anwendung der Vandermonde-Faltung erh\"alt man
$R(k,k) \le \binom{2k-2}{k-1}$.  Die Schranke $\binom{2k-2}{k-1} \le 4^{k-1}$
folgt aus $\binom{2m}{m} \le 4^m$.
\end{proof}

\subsection{Die Spencer-Verfeinerung (1975)}

Spencer \cite{Spencer1975} verbesserte die Erd\H{o}s-Schranke mit Hilfe des
Lov\'asz-Lokalen-Lemmas (vgl.\ Abschnitt~\ref{sec:lll}):

\begin{theorem}[Spencer 1975]\label{satz:spencer}
Es existiert eine absolute Konstante $c > 0$ mit
\[
  R(k,k) \;\ge\; c \cdot k \cdot \frac{2^{k/2}}{\sqrt{\ln k}}.
\]
\end{theorem}

Diese Verbesserung gegen\"uber der einfachen Erd\H{o}s-Schranke betr\"agt einen
Faktor der Ordnung $k/\sqrt{\ln k}$.

\subsection{Die Sah-Schranke (2023)}

Ein wichtiger j\"ungerer Durchbruch stammt von Sah \cite{Sah2023}:

\begin{theorem}[Sah 2023]
Es existiert $\varepsilon > 0$ (explizit: $\varepsilon \approx 0{,}007$) mit
\[
  R(k,k) \;\le\; (4-\varepsilon)^k.
\]
\end{theorem}

Dies ist die erste Verbesserung der $4^k$-Oberschranke seit \"uber 75 Jahren
und durchbricht die Erd\H{o}s-Szekeres-Schranke.

%% ============================================================
\section{Der Hales--Jewett-Satz}
%% ============================================================

Der Satz von Hales und Jewett \cite{HalesJewett1963} ist ein kombinatorischer
Grundstein, der dem Satz von van der Waerden \"uber arithmetische Progressionen
zugrunde liegt.

\begin{definition}
Eine \emph{kombinatorische Linie} im W\"urfel $[k]^n = \{1,\ldots,k\}^n$ ist eine
Menge von $k$ Punkten $\ell(1), \ell(2), \ldots, \ell(k)$, bei der jede Koordinate
$i$ entweder f\"ur alle $j$ gleich ist (eine \emph{feste} Koordinate) oder
$\ell(j)_i = j$ erf\"ullt (eine \emph{aktive} Koordinate), mit mindestens einer
aktiven Koordinate.
\end{definition}

\begin{theorem}[Hales--Jewett, 1963]\label{satz:hales-jewett}
F\"ur alle $k \ge 1$ und $r \ge 1$ existiert $N = \mathrm{HJ}(k,r)$, sodass
jede $r$-F\"arbung von $[k]^N$ eine einfarbige kombinatorische Linie enth\"alt.
\end{theorem}

\begin{proof}[Beweissk\"otte]
Der Beweis erfolgt durch doppelte Induktion \"uber $(k,r)$.  Die Basisf\"alle
sind trivial.  Im Induktionsschritt bettet man $[k]^n$ in einen gr\"o\ss{}eren
W\"urfel ein und verwendet die Induktionsvoraussetzung, um eine einfarbige
kombinatorische Linie durch ein Dichteargument zu finden.  Der Dichte-Hales--Jewett-Satz
(Polymath 2012) liefert quantitative Schranken.
\end{proof}

\begin{corollary}[Satz von van der Waerden]
F\"ur alle $k,r \ge 1$ existiert $W = W(k,r)$, sodass jede $r$-F\"arbung von
$\{1,\ldots,W\}$ eine einfarbige arithmetische Progression der L\"ange $k$ enth\"alt.
\end{corollary}

\begin{proof}
Das Korollar folgt aus dem Hales--Jewett-Satz durch Projektion der kombinatorischen
Linie in $[k]^N$ via der Abbildung $(x_1,\ldots,x_N) \mapsto \sum_{i=1}^N k^{i-1}(x_i-1)$:
Kombinatorische Linien projizieren auf arithmetische Progressionen.
\end{proof}

%% ============================================================
\section{Extremale Kombinatorik}
%% ============================================================

\subsection{Der Tur\'an-Satz}

Die extremale Graphentheorie fragt: Was ist die maximale Kantenanzahl eines
Graphen auf $n$ Knoten, der $K_{r+1}$ nicht als Teilgraphen enth\"alt?

\begin{definition}
Die \emph{Tur\'an-Zahl} $\mathrm{ex}(n, K_{r+1})$ ist die maximale Kantenanzahl
in einem $K_{r+1}$-freien Graphen auf $n$ Knoten.  Der \emph{Tur\'an-Graph}
$T(n,r)$ ist der vollst\"andige $r$-partite Graph auf $n$ Knoten mit m\"oglichst
gleich gro\ss{}en Teilen.
\end{definition}

\begin{theorem}[Tur\'an, 1941]\label{satz:turan}
\[
  \mathrm{ex}(n, K_{r+1}) \;=\; |E(T(n,r))| \;=\; \left(1 - \frac{1}{r}\right)\frac{n^2}{2}
  + O(n).
\]
Genauer gilt $\mathrm{ex}(n, K_{r+1}) = \left\lfloor\left(1-\frac{1}{r}\right)\frac{n^2}{2}\right\rfloor$
und der eindeutige Extremalgraph ist $T(n,r)$.
\end{theorem}

\begin{proof}
\textbf{Oberschranke.}  Sei $G$ ein $K_{r+1}$-freier Graph auf $n$ Knoten.
Durch Induktion \"uber $n$: Hat $G$ einen Knoten $v$ mit Maximalgrad $\Delta$,
so sei $G' = G - v$.  Nach Induktionsvoraussetzung ist
$|E(G')| \le \left(1-\frac{1}{r}\right)\frac{(n-1)^2}{2}$.
Da $G$ $K_{r+1}$-frei ist, ist die Nachbarschaft $N(v)$ $K_r$-frei; eine sorgf\"altige
Rechnung liefert die Schranke.

\textbf{Unterschranke.}  Der Tur\'an-Graph $T(n,r)$ ist $K_{r+1}$-frei und erreicht
die angegebene Kantenanzahl.

\textbf{Eindeutigkeit.}  Ein Stabilit\"atsargument zeigt, dass jeder Extremalgraph
$r$-f\"arbbar und daher gleich $T(n,r)$ sein muss.
\end{proof}

\subsection{Das Zarankiewicz-Problem}

\begin{definition}
$z(n;r)$ bezeichnet die maximale Kantenanzahl in einem bipartiten Graphen auf
$n+n$ Knoten, der keinen vollst\"andigen bipartiten Teilgraphen $K_{r,r}$ enth\"alt.
\end{definition}

\begin{theorem}[K\H{o}v\'ari--S\'os--Tur\'an, 1954]
\[
  z(n;r) \;\le\; \frac{1}{2}\left(1+\sqrt{4r-3}\right)(n-1)^{1-1/r}\,n^{1/r}
  + \frac{r-1}{2}\,n \;=\; O\!\left(n^{2-1/r}\right).
\]
\end{theorem}

%% ============================================================
\section{Die Probabilistische Methode und das Lov\'asz-Lokale-Lemma}
\label{sec:lll}
%% ============================================================

\subsection{Die Probabilistische Methode}

Die Probabilistische Methode beweist die Existenz kombinatorischer Strukturen,
indem gezeigt wird, dass ein zuf\"allig gew\"ahltes Objekt die gew\"unschte
Eigenschaft mit positiver Wahrscheinlichkeit hat.  Wichtige Werkzeuge sind:
\begin{itemize}[noitemsep]
  \item \textbf{Erstes-Moment-Methode} (Vereinigungsschranke):
        $\Pr[\bigcup A_i] \le \sum \Pr[A_i]$.
  \item \textbf{Zweites-Moment-Methode}: Nutzt die Varianz um
        $\Pr[X > 0] > 0$ zu zeigen.
  \item \textbf{Alteration}: Modifiziere eine zuf\"allige Struktur, um Defekte zu beheben.
\end{itemize}

\subsection{Das Lov\'asz-Lokale-Lemma}

Das Lov\'asz-Lokale-Lemma \cite{ErdosLovasz1975} ist ein m\"achtiges Werkzeug
f\"ur den Umgang mit nahezu unabh\"angigen schlechten Ereignissen.

\begin{theorem}[Lov\'asz-Lokales-Lemma -- Symmetrische Version]\label{satz:lll}
Seien $A_1, \ldots, A_n$ Ereignisse in einem Wahrscheinlichkeitsraum, jedes mit
$\Pr[A_i] \le p$, und jedes $A_i$ sei von allen bis auf h\"ochstens $d$ anderen
Ereignissen stochastisch unabh\"angig.  Gilt
\[
  e\,p\,(d+1) \;\le\; 1,
\]
so ist $\Pr\bigl[\bigcap_{i=1}^n \overline{A_i}\bigr] > 0$.  Insbesondere
existiert ein Elementarereignis, bei dem kein $A_i$ eintritt.
\end{theorem}

\begin{proof}
Wir beweisen die allgemeine (asymmetrische) Version und leiten den symmetrischen
Fall her.  Definiere einen \emph{Abh\"angigkeitsgraphen} $D$, in dem $A_i$ und
$A_j$ verbunden sind, wenn sie nicht unabh\"angig sind.  Seien $x_i \in [0,1)$
mit $x_i = ep$.  Durch Induktion \"uber $|S|$ zeigt man:
\[
  \Pr\left[A_i \;\middle|\; \bigcap_{j \in S} \overline{A_j}\right]
  \;\le\; x_i \prod_{j \sim i} (1-x_j).
\]
F\"ur $ep(d+1) \le 1$ gilt $x_i \prod_{j\sim i}(1-x_j) \ge ep\cdot(1-ep)^d \ge \Pr[A_i]$.
Daraus folgt
$\Pr\bigl[\bigcap \overline{A_i}\bigr] \ge \prod_{i=1}^n (1-x_i) > 0$.
\end{proof}

\subsection{Anwendung: Graphenf\"arbung}

\begin{corollary}
Hat der Graph $G$ Maximalgrad $\Delta \ge 3$ und Taillenweite $g \ge 5$, so gilt
f\"ur die chromatische Zahl
\[
  \chi(G) \;\le\; \left\lceil \frac{2\Delta}{\ln \Delta} \right\rceil + 1.
\]
\end{corollary}

\subsection{Anwendung auf Ramsey-Zahlen}

Mit dem LLL bewies Spencer \cite{Spencer1975}:

\begin{theorem}
F\"ur $k \ge 3$ existiert eine Zweif\"arbung von $K_n$ ohne einfarbiges $K_k$
f\"ur
\[
  n \;\ge\; \frac{k-1}{e\sqrt{2}} \cdot \frac{2^{k/2}}{\sqrt{\ln k}}.
\]
Daher gilt $R(k,k) > c \cdot k \cdot 2^{k/2}/\sqrt{\ln k}$.
\end{theorem}

%% ============================================================
\section{Offene Probleme}
%% ============================================================

\subsection{Der Wert von \texorpdfstring{$R(5,5)$}{R(5,5)}}

\begin{conjecture}[Exakter Wert von $R(5,5)$]\label{verm:r55}
Die Ramsey-Zahl $R(5,5)$ erf\"ullt
\[
  43 \;\le\; R(5,5) \;\le\; 48.
\]
Der exakte Wert ist unbekannt.
\end{conjecture}

Die Untergrenze $R(5,5) \ge 43$ wurde von Exoo \cite{Exoo1989} durch einen
expliziten Graphen auf 42 Knoten ohne einfarbiges $K_5$ etabliert.  Die Obergrenze
$R(5,5) \le 48$ wurde von Angeltveit und McKay \cite{AngeltveitMcKay2017} bewiesen.
Die Bestimmung des exakten Wertes ist eines der bekanntesten offenen Probleme der
Kombinatorik.

\subsection{Diagonale Ramsey-Zahlen}

\begin{conjecture}[Erd\H{o}s]\label{verm:diagonal}
Es existieren Konstanten $c_1 < c_2$ mit
\[
  c_1^k \;\le\; R(k,k) \;\le\; c_2^k
\]
f\"ur alle $k \ge 2$.  Dar\"uber hinaus existiert der Grenzwert
$\lim_{k\to\infty} R(k,k)^{1/k}$.
\end{conjecture}

Trotz der Schranken $2^{k/2} < R(k,k) \le (4-\varepsilon)^k$ ist die Frage,
ob dieser Grenzwert existiert und welchen Wert er hat, ungekl\"art.

\subsection{Off-diagonale Ramsey-Zahlen}

\begin{conjecture}[Erd\H{o}s-Faudree]\label{verm:offdiag}
F\"ur festes $s \ge 3$ gilt
\[
  R(s,t) \;=\; \Theta\!\left(t^{(s+1)/2} / (\log t)^{(s-1)/2}\right).
\]
\end{conjecture}

Der Fall $s=3$: $R(3,t) = \Theta(t^2/\log t)$ wurde von Kim \cite{Kim1995} bewiesen.

\subsection{Mehrfarben-Ramsey-Zahlen}

\begin{conjecture}
Die Mehrfarben-Ramsey-Zahl $R_r(3)$ (kleinste $N$, sodass jede $r$-F\"arbung von
$K_N$ ein einfarbiges $K_3$ enth\"alt) erf\"ullt
\[
  R_r(3) \;=\; \Theta\!\left(\frac{r!}{\sqrt{r}}\right).
\]
\end{conjecture}

\subsection{Erd\H{o}s-Preisaufgaben}

Erd\H{o}s bot Geldpreise f\"ur die L\"osung vieler Probleme dieses Bereichs an:
\begin{itemize}[noitemsep]
  \item \$250 f\"ur die Bestimmung von $\lim_{k} R(k,k)^{1/k}$,
  \item \$500 f\"ur den Beweis von $R(3,t) = o(t^2/\log t)$ oder $\omega(t^2/\log t)$,
  \item \$10.000 (h\"ochster Preis) f\"ur eine L\"osung der Frage nach den Ramsey-Multiplikationskonstanten.
\end{itemize}

%% ============================================================
\section*{Literatur}
%% ============================================================

\begin{thebibliography}{99}

\bibitem{Ramsey1930}
F.P.~Ramsey,
\textit{On a problem of formal logic},
Proc.\ London Math.\ Soc.\ \textbf{30} (1930), 264--286.

\bibitem{Erdos1947}
P.~Erd\H{o}s,
\textit{Some remarks on the theory of graphs},
Bull.\ Amer.\ Math.\ Soc.\ \textbf{53} (1947), 292--294.

\bibitem{ErdosSzekeres1935}
P.~Erd\H{o}s und G.~Szekeres,
\textit{A combinatorial problem in geometry},
Compositio Math.\ \textbf{2} (1935), 463--470.

\bibitem{ErdosLovasz1975}
P.~Erd\H{o}s und L.~Lov\'asz,
\textit{Problems and results on 3-chromatic hypergraphs and some related questions},
in: \textit{Infinite and Finite Sets}, Colloq.\ Math.\ Soc.\ J.\ Bolyai \textbf{10},
North-Holland, 1975, S.~609--627.

\bibitem{GrahamRothschildSpencer}
R.L.~Graham, B.L.~Rothschild und J.H.~Spencer,
\textit{Ramsey Theory}, 2.~Aufl., Wiley, New York, 1990.

\bibitem{AlonSpencer}
N.~Alon und J.H.~Spencer,
\textit{The Probabilistic Method}, 4.~Aufl., Wiley, Hoboken, 2016.

\bibitem{HalesJewett1963}
A.W.~Hales und R.I.~Jewett,
\textit{Regularity and positional games},
Trans.\ Amer.\ Math.\ Soc.\ \textbf{106} (1963), 222--229.

\bibitem{Spencer1975}
J.H.~Spencer,
\textit{Ramsey's theorem for spaces},
Trans.\ Amer.\ Math.\ Soc.\ \textbf{249} (1975), 187--198.

\bibitem{Turan1941}
P.~Tur\'an,
\textit{On an extremal problem in graph theory} (ungarisch),
Mat.\ Fiz.\ Lapok \textbf{48} (1941), 436--452.

\bibitem{Exoo1989}
G.~Exoo,
\textit{A lower bound for Schur numbers and multicolor Ramsey numbers of $K_3$},
Electron.\ J.\ Combin.\ \textbf{1} (1994), R8.

\bibitem{AngeltveitMcKay2017}
J.~Angeltveit und B.D.~McKay,
\textit{$R(5,5) \le 48$},
J.\ Graph Theory \textbf{89} (2018), Nr.~1, 5--32.

\bibitem{Kim1995}
J.H.~Kim,
\textit{The Ramsey number $R(3,t)$ has order of magnitude $t^2/\log t$},
Random Structures Algorithms \textbf{7} (1995), 173--207.

\bibitem{Sah2023}
A.~Sah,
\textit{Diagonal Ramsey via effective quasirandomness},
Duke Math.\ J.\ \textbf{172} (2023), Nr.~3, 545--567.

\bibitem{Radziszowski2021}
S.P.~Radziszowski,
\textit{Small Ramsey numbers},
Electron.\ J.\ Combin., Dynamic Survey DS1, Revision 16, 2021.

\end{thebibliography}

\end{document}
